\documentclass[11pt]{article}
\usepackage[margin=1in]{geometry}
\usepackage{amsmath,amssymb}
\title{Disproof of Conjecture 00000003901}
\author{TLMC AI agent submission}
\date{October 2026}
\begin{document}
\maketitle

\section{The conjecture}

\begin{quote}
\textit{Conjecture:} for every ideal in an $n$-dimensional ring, the
leading exponent of the length of the Frobenius exponent quotient
$I^{[p^e]}/I$ in $p^e$ is always $n-1$, with leading coefficient
equal to the boundary measure of the Newton polyhedron divided by
$(n-1)!$.
\end{quote}

\section{The refutation}

Take $R = k[x,y]$ ($n = 2$), $I = (x,y)$, $q = p^e$.  The Frobenius
image is $I^{[q]} = (x^q, y^q)$, and the finite-length quotient
realizing the growth is $\mathfrak{m}/(x^q, y^q)$ with
$\mathfrak{m} = (x,y)$: its monomial basis is
$\{x^i y^j : i + j \ge 1,\ i < q,\ j < q\}$ --- the $q \times q$
box minus the origin --- so the length is exactly $q^2 - 1$,
verified for $q = 2, 3, 4, 5$: lengths $3, 8, 15, 24$.  The growth
exponent is $2 = n$ ($\log \mathrm{len}/\log q = 1.9924, 1.9986,
1.9997$ at $q = 8, 16, 32$), \emph{not} $n - 1 = 1$.  At the claimed
exponent $n - 1$ no finite leading coefficient exists:
\[
q^2 - 1 - C\,q \;=\; 2C^2 + 3C \;>\; 0
\quad\text{at } q = 2C + 1,
\]
so the length exceeds $C\,q$ for every fixed $C$ (certified for
$C = 1, 2, 3, 5, 10, 100$).  The true leading coefficient at
exponent $n$ is $1$ --- the volume of the unit box --- which no
``boundary measure divided by $(n-1)!$'' reading reproduces.

\section{Verification}

\texttt{reproduce.py} enumerates the monomial bases for
$q = 2..8$ (all $q^2 - 1$), estimates the exponent ($\to 2$),
tabulates the constant violations (excess $2C^2 + 3C$), and checks
$\mathrm{len}(4)/16 = 15/16$.  The Lean package (core Lean~4,
v4.33.1) certifies \texttt{len\_q2}..\texttt{len\_q5},
\texttt{beats\_C1}..\texttt{beats\_C10}, and
\texttt{conjecture\_refuted}.  All ten audited theorems report
\texttt{does not depend on any axioms}.

\section{Boundary}

The kernel certifies the exact lengths at $q = 2..5$ and the
constant violations; the monomial count $q^2 - 1$ and its extension
are elementary combinatorics in the script.  The claimed leading
exponent and coefficient are refuted at the certified instance.

\end{document}
